\documentclass[11pt]{article}
\usepackage[T1]{fontenc}
\usepackage{amsmath}
\usepackage{amssymb}
\usepackage{amsthm}
\usepackage[hidelinks]{hyperref}
\newtheorem{theorem}{Theorem}[section]
\newtheorem{lemma}[theorem]{Lemma}
\newtheorem{corollary}[theorem]{Corollary}
\theoremstyle{definition}
\newtheorem{definition}[theorem]{Definition}
\begin{document}
\begin{center}
{\LARGE Natural number order}
\end{center}
\section{Positivity}
\label{ll:order:positivity}
\begin{definition}
\label{ll:order:positive}
For every natural number \(n\), \(n\) is positive holds exactly when \(0 < n\).
\end{definition}
\begin{theorem}
\label{ll:order:successor-positive}
For every natural number \(n\), the successor of \(n\) is positive.
\end{theorem}
\begin{proof}
The goal follows from \(\operatorname{zeroltsucc}(n)\).
\end{proof}
\section{Order}
\label{ll:order:induction}
\begin{theorem}
\label{ll:order:zero-le}
For every natural number \(n\), \(0 \le n\).
\end{theorem}
\begin{proof}
Proceed by induction on \(n\).
Case zero:
The goal follows from \(\operatorname{lerefl}(0)\).
Case \(\operatorname{succ}\) with \(m\), \(ih\):
The goal follows from \(\operatorname{letrans}(ih, \operatorname{lesucc}(m))\).
\end{proof}
\begin{theorem}
\label{ll:order:le-add-right}
For every natural number \(n\), \(n \le n + 0\).
\end{theorem}
\begin{proof}
Rewrite the goal using \(\rightarrow \texttt{Addition::add-zero}(n)\).
The goal follows from \(\operatorname{lerefl}(n)\).
\end{proof}
\section{Case analysis}
\label{ll:order:case-analysis}
\begin{theorem}
\label{ll:order:zero-or-successor}
For every natural number \(n\), \(n = 0\) or there exists a natural number \(m\) such that \(n = \operatorname{succ}(m)\).
\end{theorem}
\begin{proof}
Consider the cases of \(n\).
Case zero:
Select the left alternative.
The goal follows by reflexivity.
Case \(\operatorname{succ}\) with \(k\):
Select the right alternative.
Use \(k\) as the witness.
The goal follows by reflexivity.
\end{proof}
\begin{theorem}
\label{ll:order:successor-not-zero}
For every natural number \(n\), not \(\operatorname{succ}(n) = 0\).
\end{theorem}
\begin{proof}
Assume \(h\).
Apply \(\operatorname{succnezero}(n)\).
The goal follows from \(h\).
\end{proof}
\end{document}
